\subsubsection{递推顺序、支配筛选与原箱号恢复}
某区有$K_i$个等价类，需求计数为$d_{ik}$，计数状态数量为
\begin{equation}
|\mathcal N_i|=\prod_{k=1}^{K_i}(d_{ik}+1).
\label{eq:q1_dp_states}
\end{equation}
每个非空批型至少减少一箱，故按箱数递增计算状态时，前驱已经处理。原点为$(0,0,0)$，不可达状态保持空标记，不能被默认填成零成本。直接枚举状态与批型的实现时间为$O(|\mathcal N_i||\mathcal P_i|)$，单目标词典序标签及前驱需要$O(|\mathcal N_i|)$存储。

每次更新同时保存前驱状态和所选批型。最终从完整需求状态回溯，获得机型和计数批型；再从相应等价类的原始箱号池中取出对应数量，恢复唯一箱号。恢复后按原质量、体积和完整往返能耗重新核验，计数压缩不会把货箱变成可分割的连续质量。

多目标递推中只在相同剩余需求状态内比较支配关系。若标签$u$在架次、能耗、时间均不大于$v$，则它们接上相同后续批型后仍保持支配，故$v$可删除；剩余需求不同的状态不能这样互相删除。各区标签集再逐次求和并筛除支配目标。Bouman等在无人机路径研究中使用动态规划压缩重复决策\cite{bouman2018}，本文采用的状态则直接对应本题的不可拆需求计数。

\begin{table}[htbp]\centering\small
\caption{精确组批递推与方案恢复步骤}\label{tab:dp_execution}
\begin{tabularx}{\linewidth}{p{1.5cm}L}\toprule
阶段 & 操作\\\midrule
输入 & 读取等价类需求及质量、体积、能源均合格的计数批型。\\
初始化 & 原点赋零标签，其他状态不可达；设置前驱记录。\\
递推 & 按箱数递增遍历，尝试不超过当前需求的批型，累加真实成本。\\
筛选 & 词典序求解保留最优标签；多目标求解保留同状态非支配标签。\\
恢复 & 回溯机型与批型，从原箱号池中分配完整货箱。\\
核验 & 核对箱号全集、容量、返航能量及与整数模型的目标一致性。\\\bottomrule
\end{tabularx}
\end{table}
